\par\Needspace{4\baselineskip}
\begin{center}
\SourceLeaf{025b}\pdfbookmark[2]{7.}{article-007}7.
\end{center}
Im Perihel sind also die wahre, die exzentrische und die mittlere Anomalie $= 0$.
Wenn die wahre Anomalie dann zunimmt, nehmen auch die exzentrische und die mittlere zu.
Dabei bleibt die exzentrische kleiner als die wahre und die mittlere kleiner als die
exzentrische, bis im Aphel alle drei zugleich $= 180^0$ werden. Von dort bis zum Perihel
ist die exzentrische Anomalie stets größer als die wahre und die mittlere größer als
die exzentrische, bis im Perihel alle drei $= 360^0$ werden, was gleichbedeutend damit
ist, dass sie alle wieder $= 0$ sind. Entsprechen der wahren Anomalie $v$ die exzentrische $E$ und die mittlere $M$, so
ergibt sich allgemein, dass der wahren Anomalie $360^0-v$ die exzentrische
$360^0-E$ und die mittlere $360^0-M$ entsprechen.
Die Differenz $v-M$ zwischen der wahren und der mittleren Anomalie heißt
\emph{Mittelpunktsgleichung}. Sie ist somit vom Perihel bis zum Aphel positiv und
vom Aphel bis zum Perihel negativ. Im Perihel und im Aphel selbst verschwindet sie.
Da nun $v$ und $M$ den ganzen Kreis von $0$ bis $360^0$ in derselben Zeit durchlaufen,
erhält man die Zeit eines Umlaufs, die auch \emph{Umlaufzeit} heißt, in Tagen, indem
man $360^0$ durch die tägliche Bewegung $\frac{k\BareRoot{(1+\mu)}}{a^{\frac{3}{2}}}$
teilt. Daraus ergibt sich, dass bei verschiedenen Körpern, die die Sonne umlaufen,
die Quadrate der Umlaufzeiten proportional zu den Kuben der mittleren Abstände sind,
soweit sich die Massen der Körper oder, genauer gesagt, deren Unterschiede
vernachlässigen lassen.
